\subsection{A Jobs Guarantee, and the Objection It Has to Answer}
\label{sec:guarantee-objection}
Where automation displaces work, the standard response is retraining, funded by whoever can be persuaded to fund it. The design examined here is more direct: a jobs guarantee, with the state as employer of last resort, hiring anyone who applies, where they live, at a living wage, with paid on-the-job training built into the position rather than offered in place of one. It is examined for two reasons. It is a public institution that has to survive an administrator hostile to it, which is this book's custody problem in another setting; and chapter~\ref{ch:standing} places a bearer that has left its post into public work on this model, so the three conditions stated below are the ones §\ref{sec:reserve-behind} inherits. The jurisdiction in this section and the next is the United States, whose federal structure is what the design leans on.

\runin{The objection this proposal actually has to answer}
A guarantee is only as good as the state administering it, and employer of last resort under a hostile administration is a placement lever with a compliance condition attached. Refuse the placement and lose the guarantee. Speak against the government that employs you and find the placement changes. Everything this book says about how a channel becomes an instrument of the thing it was built to check applies to this proposal. One instance from 2025 makes it concrete. The Consumer Financial Protection Bureau draws its funding from the Federal Reserve on the director's request rather than by annual appropriation, an arrangement the Supreme Court had upheld the year before; in February 2025 an acting director, removable at will, declined to make the request, and the agency was stopped from inside without a statute being amended or a court being asked \autocite{cfpb2024cfsa,kolodziej2025cfpb}.

The answer is a design specification and not a reassurance, because reassurance cannot be audited. Three conditions:

\begin{itemize}
  \item Placements are appealable, to a body that did not make them, on a record that survives the appeal. A placement decision that leaves no artifact is a placement decision that cannot be reviewed.
  \item Refusal rights are real and priced by the person refusing. A worker who can decline a placement without losing the guarantee has a refusal right; one who can decline only by exiting the program has none: a party with no way out has voice without exit (§\ref{sec:exit-leverage}), and can object and be overruled.
  \item The people who audit placements are not the people who make them.
\end{itemize}

Of the shapes available, federally funded and locally administered is the one that meets the three conditions, and it meets them on redundancy rather than on local fit: many administrators, separately capturable, mean there is no single point at which a hostile executive flips the program, and capturing a thousand local authorities is a more visible undertaking than capturing one federal one. It cuts both ways, since local administration is also where a hostile local authority attaches conditions of its own, and there the federal funder is the check. One distinction has to be explicit, because the same words describe a different policy. Devolution is also how a program is shrunk: hand the obligation down without the money, and it can be starved locally while the center points at local administration. The test that separates them is checkable — whether funding is unconditional and centrally guaranteed while placement is local. If funding can be withheld from a jurisdiction, the redundancy argument is void and the structure has become the lever it was supposed to prevent.

A guarantee that can be revoked by whoever is in office is not a guarantee, and everything the next section claims for the design rests on the three conditions above.
